\section{The moderate-row seed}\label{sec:moderate}

\begin{theorem}\label{thm:moderate}
There are absolute constants $A,C_m,\eps_0>0$ such that the following holds.
Let $U\in\R^{n\times d}$ be Parseval with $2d\le n$, $d\ge A\log(2n)$, and
$|\norm{u_i}^2-a|\le\eps a$ for all $i$, where $0<\eps\le\eps_0$. Then there is
an ENP frame $\widehat U$ with $\fro{\widehat U-U}^2\le C_m\eps d$.
\end{theorem}

Throughout this section $P=UU^T$, $p=\diag P$, $D_p=\Diag(p)$, $L=L_P$, and
$\eps_0\le\frac12$, so $\frac a2\le p_i\le\frac{3a}2\le\frac34$. For a symmetric
$M\in\R^{n\times n}$, $\Lap(M)=\sum_{i<j}M_{ij}^2(e_i-e_j)(e_i-e_j)^T$ is the
Laplacian with weights $M_{ij}^2$; thus $L=\Lap(P)$ and $x^T\Lap(M)x=\frac12\fro{[\Diag
x,M]}^2$. For nonnegative symmetric weights $w$, $\Lap[w]$ denotes the
Laplacian with weights $w_{ij}$. We use $\opn{\Lap[w]}\le2\max_i\sum_jw_{ij}$.

\subsection{The horizontal tangent space and the filtered covariance}

Fix an isometry $V\in\R^{n\times(n-d)}$ onto $(\im U)^\perp$, and let
$\mathcal Z=\R^{(n-d)\times d}$ with the Frobenius inner product. For
$Z\in\mathcal Z$ the matrix $H=VZ$ satisfies $U^TH=0$; its rows are
$h_i=Z^Tv_i$. Define
\[
 A:\mathcal Z\to\R^n,\quad (AZ)_i=v_i^TZu_i=(HU^T)_{ii},\qquad
 A^*x=N_x:=V^T\Diag(x)U .
\]
Then $AA^*=L$, i.e.\ $\fro{N_x}^2=x^TLx$ (\cref{lem:energy}). Let
\[
 \bar A=D_p^{-1/2}A,\qquad S=\bar A\bar A^*=D_p^{-1/2}LD_p^{-1/2},\qquad
 B=\bar A^*\bar A .
\]
Since $L\succeq0$ and $S=I-D_p^{-1/2}(P\circ P)D_p^{-1/2}$ with $P\circ P\succeq0$
(Schur product theorem),
\begin{equation}\label{eq:S}
 0\preceq S\preceq I,\qquad 0\preceq B\preceq I,\qquad
 \tr(I-S)=\sum_i\frac{P_{ii}^2}{p_i}=d .
\end{equation}
Fix an orthonormal eigenbasis of $S$, and for each of its members $y$ with
eigenvalue $\mu>0$ put $f_y=D_p^{-1/2}y$ and
$Z_y=\mu^{-1/2}\bar A^*y=\mu^{-1/2}N_{f_y}$. Since $BZ_y=\mu Z_y$ and
$\ip{Z_y}{Z_{y'}}=\ip y{Sy'}/\sqrt{\mu\mu'}$, the $Z_y$ form an orthonormal
eigenbasis of $B$ on $(\ker B)^\perp=\im\bar A^*$; and since $p_i\ge a/2$,
\begin{equation}\label{eq:finfty}
 \norm{f_y}_\infty\le\sqrt{2/a}.
\end{equation}

\begin{definition}[Filtered covariance]
For $\rho>0$ let $C_\rho=\rho(I-B)(\rho I+B)^{-1}$ on $\mathcal Z$, and let
$Z\sim N(0,C_\rho/n)$, $H=VZ$, $Y=HU^T+UH^T$.
\end{definition}

\begin{lemma}\label{lem:filter}
\begin{enumerate}[label=(\alph*)]
\item $0\preceq C_\rho\preceq I$; $C_\rho=I$ on $\ker B$.
\item $\Cov(AZ)=n^{-1}D_p^{1/2}\,\rho S(I-S)(\rho I+S)^{-1}D_p^{1/2}\preceq
(\rho/n)D_p$.
\item $I-C_\rho=B+R_\rho$ with $R_\rho=\sum_{\mu>0}r_\mu\,Z_y\otimes Z_y$,
$r_\mu=\mu(1-\mu)/(\rho+\mu)\ge0$, and
\[
 \sum_{\mu>0}r_\mu\le d,\qquad \sum_{\mu>0}\frac{r_\mu}{\sqrt\mu}\le\frac d{2\sqrt\rho},
 \qquad\sum_{\mu>0}\frac{r_\mu}\mu\le\frac d\rho .
\]
\end{enumerate}
\end{lemma}

\begin{proof}
(a) The eigenvalues of $C_\rho$ are $\rho(1-\mu)/(\rho+\mu)\in[0,1]$. (b) By
singular value calculus $\bar AC_\rho\bar A^*=\rho S(I-S)(\rho I+S)^{-1}$, whose
eigenvalues $\rho\mu(1-\mu)/(\rho+\mu)$ are at most $\rho$. (c) On the
$\mu$-eigenspace, $1-\rho(1-\mu)/(\rho+\mu)=\mu+\mu(1-\mu)/(\rho+\mu)$. The sums
follow from $\mu/(\rho+\mu)\le1$, $\sqrt\mu/(\rho+\mu)\le1/(2\sqrt\rho)$,
$1/(\rho+\mu)\le1/\rho$, and $\sum_{\mu>0}(1-\mu)\le\tr(I-S)=d$.
\end{proof}

\subsection{The second-order mean}

For $Z\in\mathcal Z$ define the quadratic diagonal
\[
 q_i(Z)=\norm{h_i}^2-\norm{Zu_i}^2=\norm{Z^Tv_i}^2-\norm{Zu_i}^2 .
\]
For a covariance $\Sigma$ on $\mathcal Z$ write $q(\Sigma)=\E q(G)$ with
$G\sim N(0,\Sigma)$; $q(\cdot)$ is linear, and $q(\zeta\otimes\zeta)=q(\zeta)$.

\begin{lemma}[Mean]\label{lem:mean}
$\E q(Z)=a\one-p-b_0-b_*$, where
\[
 b_0=\frac1nL(p^{-1}),\quad \norm{b_0}_\infty\le\frac{12\eps}n,\qquad
 b_*=\frac1n\sum_{\mu>0}r_\mu\,q(Z_y),
\]
and $b_*\perp\one$ with
\begin{equation}\label{eq:bias-dual}
 \ip x{b_*}^2\le\frac{2a}\rho\,x^TLx\qquad(x\in\R^n).
\end{equation}
\end{lemma}

\begin{proof}
For $\Sigma=I$: $\E\norm{Z^Tv_i}^2=d\norm{v_i}^2=d(1-p_i)$ and
$\E\norm{Zu_i}^2=(n-d)p_i$, so $q(I/n)=a\one-p$. Next
$B=\sum_j\zeta_j\otimes\zeta_j$ with $\zeta_j=\bar A^*e_j=v_ju_j^T/\sqrt{p_j}$, and
\[
 q_i(\zeta_j)=\frac{(I-P)_{ij}^2p_j-(1-p_j)P_{ij}^2}{p_j},\qquad
 \sum_jq_i(\zeta_j)=1-p_i-\sum_j\frac{P_{ij}^2}{p_j}+p_i=(Lp^{-1})_i ,
\]
using $\sum_j(I-P)_{ij}^2=1-p_i$ and $\sum_jP_{ij}^2=p_i$. Also
$(Lp^{-1})_i=\sum_jP_{ij}^2(p_i^{-1}-p_j^{-1})$ and
$|p_i^{-1}-p_j^{-1}|\le2\eps a/((1-\eps)a)^2$, so
$|(Lp^{-1})_i|\le(1+\eps)a\cdot8\eps/a\le12\eps$. With \cref{lem:filter}(c)
this gives the formula for $\E q(Z)=q(C_\rho/n)$.

For~\eqref{eq:bias-dual} put $X=\Diag(x)$, $D_x=V^TXV$, $A_x=U^TXU$. Then
$\ip x{q(Z)}=\ip Z{D_xZ-ZA_x}$, and since $X$ and $F=\Diag(f)$ commute,
\begin{equation}\label{eq:commute}
 D_xN_f-N_fA_x=V^T[X(I-P)F-FPX]U=V^T[F(I-P)X-XPF]U=D_fN_x-N_xA_f .
\end{equation}
For $Z=Z_y=\mu^{-1/2}N_f$, $f=f_y$, this gives
$\ip x{q(Z_y)}=\mu^{-1}\ip{N_f}{D_fN_x-N_xA_f}$, and as $\fro{N_f}^2=\mu$,
$\fro{N_x}^2=x^TLx$, $\opn{D_f},\opn{A_f}\le\norm f_\infty$,
\[
 |\ip x{q(Z_y)}|\le\frac{2\norm{f_y}_\infty}{\sqrt\mu}\sqrt{x^TLx}
 \le2\sqrt{\frac2{a\mu}}\sqrt{x^TLx} .
\]
Summing with weights $r_\mu/n$ and using \cref{lem:filter}(c),
$|\ip x{b_*}|\le\frac1n\cdot\frac d{2\sqrt\rho}\cdot2\sqrt{2/a}\sqrt{x^TLx}
=\sqrt{2a/\rho}\sqrt{x^TLx}$. Taking $x=\one$ gives $b_*\perp\one$.
\end{proof}

\begin{remark}
The bias $b_*$ need not be small in $\ell_\infty$: near a direct sum whose
components are joined through a few rows it is concentrated there. What
\eqref{eq:bias-dual} says is that it can be removed by a scaling of
energy $O(at^4/\rho)$ (\cref{lem:local}), which is negligible once
$\rho\asymp t^2$ with a large constant.
\end{remark}

\subsection{The expected graph}

For a direction $\zeta\in\mathcal Z$ write
$Y_\zeta=V\zeta U^T+U\zeta^TV^T$; then $\fro{Y_\zeta}^2=2\fro\zeta^2$, and for a
covariance $\Sigma=\sum_k\lambda_k\psi_k\otimes\psi_k$,
$\E_\Sigma\Lap(Y)=\sum_k\lambda_k\Lap(Y_{\psi_k})$.

\begin{lemma}\label{lem:expected-graph}
On $\one^\perp$,
\[
 \E\Lap(Y)\succeq\frac a4I-\Bigl(\frac2n+\frac8d+\frac8\rho\Bigr)L .
\]
\end{lemma}

\begin{proof}
\emph{Unfiltered noise.} For $Z_0\sim N(0,I/n)$ and $i\ne j$,
$Y_{ij}=v_i^TZ_0u_j+v_j^TZ_0u_i$ has variance
$n^{-1}\fro{v_iu_j^T+v_ju_i^T}^2=n^{-1}(p_i+p_j-2p_ip_j-2P_{ij}^2)$. As
$p\le\frac34$, $p_i+p_j-2p_ip_j\ge\frac14(p_i+p_j)\ge\frac a4$, so
$\E\Lap(Y_{Z_0})\succeq\frac a{4n}\Lap[\one\one^T]-\frac2nL=\frac a4\Pi-\frac2nL$.

\emph{A commutator bound.} For $f\in\R^n$, $F=\Diag(f)$, put
$T_f=Y_{N_f}=(I-P)FP+PF(I-P)=FP+PF-2PFP$. Expanding,
\begin{equation}\label{eq:Tf}
 [X,T_f]=(I-2P)F[X,P]+[X,P]F(I-2P),
\end{equation}
with $I-2P$ orthogonal. Hence $x^T\Lap(T_f)x\le4\norm f_\infty^2\,x^TLx$, and for
$f=e_k$, since $E_k[X,P]$ and $[X,P]E_k$ are the $k$th row and column of $[X,P]$,
$x^T\Lap(T_{e_k})x\le2\sum_jP_{kj}^2(x_k-x_j)^2$.

\emph{Base loss.} $B=\sum_k\zeta_k\otimes\zeta_k$ with $\zeta_k=\bar
A^*e_k=N_{e_k}/\sqrt{p_k}$, so $Y_{\zeta_k}=T_{e_k}/\sqrt{p_k}$ and
$\sum_kx^T\Lap(Y_{\zeta_k})x\le\frac2a\cdot2\sum_{k,j}P_{kj}^2(x_k-x_j)^2=\frac8a\,x^TLx$.
So the base part $B/n$ of the removed covariance removes at most
$\frac8{an}L=\frac8dL$.

\emph{Residual loss.} By~\eqref{eq:finfty}, $\Lap(Y_{Z_y})=\mu^{-1}\Lap(T_{f_y})
\preceq\frac8{a\mu}L$, and by \cref{lem:filter}(c) the residual part $R_\rho/n$
removes at most $\frac8{an}\sum_\mu\frac{r_\mu}\mu L\le\frac8\rho L$.
\end{proof}

\subsection{A good sample}

\begin{lemma}\label{lem:sample}
There are absolute $h\in(0,1)$ and $b_{\max}\in\mathbb N$ and, for each
$\eta'\in(0,\frac1{16}]$, a constant $A_{\eta'}$ such that the following holds
for all $\rho>0$, $0<t\le1$, whenever $d\ge A_{\eta'}\log(2n)$. Put
$J=L+at^2\Pi$. With positive probability:
\begin{enumerate}[label=\textup{(S\arabic*)}]
\item\label{S1} $\opn Z\le16$, $\max_i\norm{h_i}^2\le3a$, $\max_i\norm{Y_i}^2\le800a$;
\item\label{S2} $\max_i|(AZ)_i|\le3\sqrt{\rho a\log(2n)/n}$ and
$\max_i|q_i(Z)-\E q_i(Z)|\le\eta'a$;
\item\label{S3} $\opn{\Lap(Y)-\E\Lap(Y)}\le\eta'a$;
\item\label{S4} $|x^T\mathcal C(Y)x|\le\frac{\eta'}t\,x^TJx$ for $x\perp\one$, where
$\mathcal C(Y)=\sum_{i<j}2P_{ij}Y_{ij}(e_i-e_j)(e_i-e_j)^T$;
\item\label{S5} there is a deterministic set $\mathfrak B$ (depending only on
$U,\rho$) with $|\mathfrak B|\le b_{\max}$ such that every $i\notin\mathfrak B$
has at least $0.95n$ indices $j\ne i$ with $|P_{ij}+tY_{ij}|\ge h\,t\sqrt{a/n}$.
\end{enumerate}
\end{lemma}

\begin{proof}
We show that each item fails with probability at most $\frac1{10}$. Note
$C_\rho\preceq I$, so every linear functional $\ip Z\Psi$ has variance at most
$\fro\Psi^2/n$; we use the tools of \cref{app:gauss}.

\ref{S1}: The net bound (\cref{lem:gauss}(d)) with $\sigma^2=1/n$ gives
$\Prob(\opn Z>16)\le2\cdot5^ne^{-8n}$. The vector $h_i=Z^Tv_i$ has covariance
$\preceq n^{-1}I_d$ and trace $\le a$; \cref{lem:gauss}(b) with $u=2a$ (so
$\min\{u^2/\fro\cdot^2,u/\opn\cdot\}\ge\min\{4d,2d\}$) and a union bound give
$\max\norm{h_i}^2\le3a$ off probability $2ne^{-cd}$. Then
$\norm{Y_i}\le\norm{h_i}+\norm{Zu_i}\le\sqrt{3a}+16\sqrt{3a/2}$.

\ref{S2}: $(AZ)_i$ is Gaussian with variance $\le\rho p_i/n\le3\rho a/(2n)$
(\cref{lem:filter}(b)); a union bound gives the first claim off probability
$2n(2n)^{-3}$. The
form $q_i(Z)=\ip Z{\mathsf M_iZ}$, $\mathsf M_iZ=v_iv_i^TZ-Zu_iu_i^T$, has
$\opn{\mathsf M_i}\le1$ and $\fro{\mathsf M_i}^2\le2(d+np_i^2)\le7d$; after
the covariance factor $(C_\rho/n)^{1/2}$ these become $\le1/n$ and
$\le\sqrt{7d}/n$, and \cref{lem:gauss}(b) with $u=\eta'a$ gives failure
probability $2n\exp(-c\eta'^2d)$.

\ref{S3}: $\Cov(\mathrm{vec}\,Y)\preceq\frac2nI$ (the map $Z\mapsto Y_Z$ is
$\sqrt2$ times an isometry) and $\E\norm{Y_i}^2\le3a$ (for unfiltered noise
$\E\norm{Y_i}^2=a(1-p_i)+(1-a)p_i$, and filtering only decreases it). Apply
\cref{lem:sqgauss}.

\ref{S4}: Let $w_{ij}=J^{-1/2}(e_i-e_j)$, so $\norm{w_{ij}}^2\le2/(at^2)$. The
matrix $t J^{-1/2}\mathcal C(Y)J^{-1/2}=\sum_{i<j}tY_{ij}\cdot2P_{ij}w_{ij}w_{ij}^T$
is a Gaussian series whose coefficients have covariance $\preceq\frac{2t^2}nI$, so
by covariance domination (\cref{lem:gauss}(a)) its matrix variance is at most
\[
 \frac{2t^2}n\sum_{i<j}4P_{ij}^2\norm{w_{ij}}^2w_{ij}w_{ij}^T
 \preceq\frac{16}{an}J^{-1/2}LJ^{-1/2}\preceq\frac{16}dI ,
\]
and \cref{lem:gauss}(a) bounds its norm by $C\sqrt{\log(20n)/d}\le\eta'$ off
probability $\frac1{10}$.

\ref{S5}: \emph{Variance bookkeeping.} By the proof of
\cref{lem:expected-graph}, for $i\ne k$ with $P_{ik}^2\le a/16$ the unfiltered
variance of $Y_{ik}$ is at least $a/(8n)$; at most $24$ indices $k$ have
$P_{ik}^2>a/16$. Row $i$ of $T_{e_j}$ is
$\delta_{ij}e_j^TP+P_{ij}e_j^T-2P_{ij}e_j^TP$, of squared norm at most
$3(\delta_{ij}p_j+P_{ij}^2+4P_{ij}^2p_j)$, so the base part removes from row $i$
total variance $\frac1n\sum_j\norm{e_i^TT_{e_j}}^2/p_j\le\frac{6}{an}\cdot5p_i\le
\frac{45}n$, which exceeds $\frac a{32n}$ at no more than $1440n/d$ positions.
The residual part removes total variance
$\frac1n\sum_\mu r_\mu\fro{Y_{Z_y}}^2\le2a$ from the whole matrix. Fix
$\delta_0=1/6400$ and let $\mathfrak B$ be the set of rows from which it removes
more than $\delta_0a$; then $|\mathfrak B|\le2/\delta_0=:b_{\max}$, and a row
$i\notin\mathfrak B$ loses more than $\frac a{32n}$ to it at no more than
$32\delta_0n$ positions. So for $d\ge10^6$ every $i\notin\mathfrak B$ has at least
$0.99n$ indices $k\ne i$ with $\Var(Y_{ik})\ge\frac a{16n}$.

\emph{Small ball.} Let $\phi:\R\to[0,1]$ be $C^1$, equal to $1$ on $[-h,h]$,
vanishing outside $[-2h,2h]$, with $|\phi'|\le2/h$, and for $i\notin\mathfrak B$ let
$F_i=\frac1n\sum_{k\ne i}\phi\bigl((P_{ik}+tY_{ik})/(t\sqrt{a/n})\bigr)$. By the
density bound for a Gaussian of variance $\ge\frac{at^2}{16n}$,
$\E F_i\le0.01+\frac{4h}{\sqrt{2\pi/16}}\le0.01+7h$. As a function of the row
$(Y_{ik})_k$, $F_i$ has gradient norm at most
$\frac1n\cdot\frac2h\cdot\sqrt{n/a}\cdot\sqrt n=2/(h\sqrt a)$; the row depends on
$Z$ through a linear map of norm $\le2$, and $Z=(C_\rho/n)^{1/2}g$ with $g$
standard. So $F_i$ is $4/(h\sqrt d)$-Lipschitz in $g$, and \cref{lem:gauss}(c) gives
$\Prob(F_i>\E F_i+0.01)\le2\exp(-c\,h^2d)$. Take $h=1/500$, so that
$\E F_i+0.01\le0.04$, and a union bound over $i\notin\mathfrak B$. On the
complement, at most $0.04n$ indices $k\ne i$ have $|P_{ik}+tY_{ik}|<ht\sqrt{a/n}$,
so at least $n-1-0.04n\ge0.95n$ have the reverse inequality.

Finally choose $A_{\eta'}$ so large that $d\ge A_{\eta'}\log(2n)$ implies $d\ge10^6$
and all the failure probabilities above are at most $\frac1{10}$ (there are
five items, so their union has probability $<1$).
\end{proof}

\subsection{Retraction}

Fix a sample as in \cref{lem:sample} and set
\[
 G=I+t^2Z^TZ,\qquad U_t=(U+tH)G^{-1/2},\qquad P_t=U_tU_t^T .
\]
Since $U^TH=0$ and $H^TH=Z^TZ$, $(U+tH)^T(U+tH)=G$ and $U_t$ is Parseval.

\begin{lemma}\label{lem:retraction}
There is an absolute $C_E$ such that for $0<t\le1$:
\begin{enumerate}[label=(\alph*)]
\item $\fro{U_t-U}^2\le C_Et^2d$;
\item every row of $\mathsf E=P_t-P-tY$ satisfies $\sum_j\mathsf E_{ij}^2\le C_Eat^4$;
\item $\diag P_t=p+2tAZ+t^2\E q(Z)+t^2(q(Z)-\E q(Z))+\mathsf r$ with
$\norm{\mathsf r}_\infty\le C_Eat^3$.
\end{enumerate}
\end{lemma}

\begin{proof}
By \ref{S1}, $\opn{G^{-1}-I},\opn{G^{-1/2}-I}\le256t^2$, $\opn{U+tH}\le17$,
$\fro{U+tH}^2\le d(1+256t^2)$ and $\norm{u_i+th_i}\le3\sqrt a$. (a)
$\fro{U_t-U}\le\fro{U+tH}\opn{G^{-1/2}-I}+t\fro Z\le Ct\sqrt d$. (b)
$\mathsf E=(U+tH)(G^{-1}-I)(U+tH)^T+t^2HH^T$, and row $i$ has norm at most
$3\sqrt a\cdot256t^2\cdot17+t^2\sqrt{3a}\cdot16$. (c) Insert
$G^{-1}=I-t^2Z^TZ+t^4(Z^TZ)^2G^{-1}$ into
$(P_t)_{ii}=(u_i+th_i)^TG^{-1}(u_i+th_i)$ and use $u_i^Th_i=(AZ)_i$:
\[
 (P_t)_{ii}=p_i+2t(AZ)_i+t^2q_i(Z)-2t^3u_i^TZ^TZh_i-t^4\norm{Zh_i}^2
 +t^4(u_i+th_i)^T(Z^TZ)^2G^{-1}(u_i+th_i),
\]
and the last three terms are $O(at^3)$ by \ref{S1}.
\end{proof}

\begin{lemma}[Regularised seed]\label{lem:regularised}
There are absolute $t_*\in(0,1]$ and $C_J$ such that if $\eta'\le\frac1{16}$,
$\rho=D^2t^2$ with $D\ge8$, and $0<t\le t_*$, then on $\one^\perp$
\[
 \tfrac1{32}J\preceq L_{P_t}\preceq C_JJ,
\]
and every $i\notin\mathfrak B$ has at least $0.9n$ indices $j\ne i$ with
$|(P_t)_{ij}|\ge\frac h2t\sqrt{a/n}$.
\end{lemma}

\begin{proof}
Entrywise $(P+tY)_{ij}^2=P_{ij}^2+2tP_{ij}Y_{ij}+t^2Y_{ij}^2$, so
$\Lap(P+tY)=L+t\mathcal C(Y)+t^2\Lap(Y)$. On $\one^\perp$, by
\ref{S4}, \ref{S3} and \cref{lem:expected-graph},
\[
 \Lap(P+tY)\succeq L-\eta'J+t^2\Bigl(\frac a4-\eta'a\Bigr)I
 -t^2\Bigl(\frac2n+\frac8d+\frac8{D^2t^2}\Bigr)L
 \succeq\Bigl(\frac78-10t^2-\frac18\Bigr)L+\frac{at^2}8I\succeq\frac18J
\]
for $t^2\le1/40$. Also $\Lap(P+tY)\preceq2L+2t^2\Lap(Y)\preceq2L+3200at^2I$
by \ref{S1}. Now $(P_t)_{ij}^2\ge\frac12(P+tY)_{ij}^2-\mathsf E_{ij}^2$ and
$(P_t)_{ij}^2\le2(P+tY)_{ij}^2+2\mathsf E_{ij}^2$, while
$\Lap(\mathsf E)\preceq2C_Eat^4I$ by \cref{lem:retraction}(b). Hence
$L_{P_t}\succeq\frac1{16}J-2C_Eat^4I\succeq\frac1{32}J$ when $t^2\le1/(64C_E)$,
and $L_{P_t}\preceq(6400+4C_E)J$. Finally, in a row $i\notin\mathfrak B$ at most
$C_Eat^4/(\frac{h^2}4t^2\frac an)=4C_Et^2n/h^2\le0.05n$ indices have
$|\mathsf E_{ij}|\ge\frac h2t\sqrt{a/n}$ once $t^2\le h^2/(80C_E)$; combine
with \ref{S5}.
\end{proof}

\subsection{Removing the bias and balancing}

\begin{proof}[Proof of \cref{thm:moderate}]
\emph{Choice of constants.} Fix $h,b_{\max}$ (\cref{lem:sample}) and
$t_*,C_J,C_E$ (\cref{lem:retraction,lem:regularised}); these are absolute. Put
$C'=1024\,e^3C_J$ and $D=\max\{250,40\sqrt{C'}/h\}$, which fixes $\rho=D^2t^2$.
Put
\[
 C_{\rm core}=\frac{32(1+b_{\max})}{h^2}+64e\,b_{\max},\qquad
 C_\Sigma=2C_E+2C',\qquad \Theta=\max\{2,C_{\rm core},C_\Sigma\},
\]
let $K_0=20\Theta$ and $\eta'=1/(20\Theta)$, and choose $A\ge A_{\eta'}$ so large
that $d\ge A\log 2n$ implies $6D\sqrt{\log(2n)/d}\le1/(20\Theta)$ and
$b_{\max}\le n/10$. Finally choose $\eps_0\le\frac12$ so that $t=\sqrt{K_0\eps}$
satisfies $t\le t_*$, $C_Et\le1/(20\Theta)$ and $12\eps\le1/(20\Theta)$ for
$\eps\le\eps_0$. The order of choices is acyclic, and none depends on
$n,d,\eps$ or $U$. Take a sample as in \cref{lem:sample}.

\emph{Bias correction.} Apply \cref{lem:local} to $W=U_t$ with $J$ as above,
$\sigma=at^2$, $c_1=\frac1{32}$, $C_1=C_J$, $f=t^2b_*$ and $r=R=\frac14$. By
\eqref{eq:bias-dual} and $L\preceq J$, $\ip fx^2\le\frac{2at^4}\rho x^TJx$, so
$E=2at^2/D^2$; with $\gamma=2c_1e^{-1}=1/(16e)$, the condition
$E<\gamma^2\sigma R^2=at^2/(4096e^2)$ holds as $D\ge250$. We obtain a Parseval
$U_*$ with
\[
 \diag P_*=\diag P_t+t^2b_*,\qquad
 \fro{U_*-U_t}^2,\ \fro{P_*-P_t}^2\le C'\frac{at^2}{D^2},\qquad
 L_{P_*}\succeq\frac{at^2}{32e}\Pi\ \text{on }\one^\perp,
\]
using $2C_1e^{4r}E/\gamma^2=2C_Je\cdot\frac{2at^2}{D^2}\cdot256e^2=C'at^2/D^2$. Since $\fro{P_*-P_t}^2\le C'at^2/D^2$, each
row loses at most $C'(at^2/D^2)/(\frac{h^2}{16}t^2\frac an)\le0.1n$ of the
indices of \cref{lem:regularised} when $D\ge40\sqrt{C'}/h$; so every
$i\notin\mathfrak B$ has at least $0.8n$ indices $j\ne i$ with
$(P_*)_{ij}^2\ge\frac{h^2}{16}\frac{at^2}n$.

\emph{Poisson constant.} \Cref{cor:core}(a) with $G=\mathfrak B^c$,
$b\le b_{\max}\le n/10$, $\gamma=h^2at^2/16$ and $\lambda=at^2/(32e)$ gives
$R\le64e\,b_{\max}/(at^2)$ and Poisson constant
$K\le(1+b_{\max})\frac{32}{h^2at^2}+\frac{64eb_{\max}}{at^2}\le\frac{C_{\rm core}}{at^2}$.

\emph{Diagonal error.} By \cref{lem:mean,lem:retraction}(c), the bias cancels
exactly:
\[
 \diag P_*-a\one=(1-t^2)(p-a\one)+2tAZ-t^2b_0+t^2\bigl(q(Z)-\E q(Z)\bigr)+\mathsf r .
\]
By \ref{S2} and $\rho=D^2t^2$,
\[
 \frac{\norm{\diag P_*-a\one}_\infty}{at^2}
 \le\frac1{K_0}+6D\sqrt{\frac{\log2n}d}+\frac{12\eps}d+\eta'+C_Et
 \le\frac5{20\Theta}=\frac1{4\Theta} .
\]

\emph{Seed and cost.} By \cref{lem:retraction}(a) and the bias correction,
$\fro{U_*-U}^2\le2C_Et^2d+2C'at^2/D^2\le C_\Sigma t^2d$ with
$C_\Sigma=2C_E+2C'$. So $U_*$ is a $(\Theta,t)$-seed for $U$, and
\cref{cor:seed} gives an ENP frame $\widehat U$ with
$\fro{\widehat U-U}^2\le4\Theta t^2d=4\Theta K_0\eps d$.
\end{proof}
